\documentclass[14pt, a4paper]{extarticle}
\usepackage[T1]{fontenc}
\usepackage[utf8]{inputenc}
\usepackage[russian]{babel}
\usepackage{amsmath, amssymb}
\usepackage{graphicx}
\usepackage{xcolor}
\usepackage{hyperref}
\usepackage{geometry}
\geometry{left=2cm, right=2cm, top=2cm, bottom=2cm}

\definecolor{codegreen}{rgb}{0,0.6,0}
\definecolor{codegray}{rgb}{0.5,0.5,0.5}
\definecolor{codepurple}{rgb}{0.58,0,0.82}
\definecolor{backcolour}{rgb}{0.95,0.95,0.92}

\usepackage[most]{tcolorbox}
\newtcolorbox{summarybox}{colback=green!5!white,colframe=green!75!black,fonttitle=\bfseries,title=Аннотация}
\newtcolorbox{importantbox}{colback=red!5!white,colframe=red!75!black,fonttitle=\bfseries,title=Основной вклад}
\newtcolorbox{infobox}{colback=blue!5!white,colframe=blue!75!black,fonttitle=\bfseries,title=Важное примечание}
\newtcolorbox{legalbox}{colback=orange!5!white,colframe=orange!75!black,fonttitle=\bfseries,title=Правовое уведомление и уведомление об интеллектуальной собственности}

\title{Deterministic Game Engine: Практическая реализация Pointer-Based Security и Local Data Regeneration Paradigms}
\author{Александр Суворов \\ \url{https://github.com/smartlegionlab}}
\date{2025}

\begin{document}

\maketitle

\begin{legalbox}
\textbf{Правовое уведомление и уведомление об интеллектуальной собственности:} Данный технический отчет описывает результаты исследований, архитектурные концепции и характеристики производительности. Все конкретные реализации, алгоритмы, исходный код и запатентованные методы защищены как интеллектуальная собственность и коммерческая тайна. Патентуемые реализации не раскрываются. Результаты производительности представлены для целей академической валидации. Данный документ фокусируется на теоретической валидации парадигм, а не на конкретных технических реализациях.
\end{legalbox}

\begin{summarybox}
\textbf{Аннотация:} В данном техническом отчете представлена валидация теоретических парадигм, изложенных в предыдущих работах автора. Исследовательский прототип, реализованный как Deterministic Game Engine, служащий модельной средой, демонстрирует архитектурные принципы, обеспечивающие генерацию бесконечных миров, симуляцию массовых NPC и верификацию состояния без передачи данных. Экспериментальные результаты предоставляют конкретные доказательства, подтверждающие теоретические преимущества предложенных парадигм, включая время доступа к состоянию, не зависящее от индекса позиции, и архитектурные паттерны без серверов.
\end{summarybox}

\textbf{Ключевые слова:} детерминированные вычисления, архитектура игрового движка, процедурная генерация, Pointer-Based Security, Local Data Regeneration, верификация, академические исследования

\begin{importantbox}
Данное исследование предоставляет экспериментальную валидацию теоретических парадигм, предложенных в предыдущих работах. Вклад демонстрирует осуществимость перехода от парадигм передачи данных к парадигмам регенерации данных, показывая потенциал для систем с улучшенными характеристиками производительности и безопасности. Результаты основаны на измерениях исследовательского прототипа.
\end{importantbox}

\begin{infobox}
\textbf{Раскрытие информации об исследовании:} В данном документе описаны архитектурные паттерны и результаты исследований. Все реализации остаются защищенной интеллектуальной собственностью. Целью является академическая валидация и установление приоритета исследований для предложенных парадигм.
\end{infobox}

\section{Введение: Контекст исследования}

Предыдущие теоретические работы \cite{suvorov2025pointer} и \cite{suvorov2025local} предложили сдвиги парадигм в подходах к безопасности данных и передаче данных. Данное исследование рассматривает вопрос практической осуществимости этих теоретических основ.

Исследовательский прототип \textbf{SMART DETERMINISTIC GAME ENGINE} предоставляет экспериментальные доказательства того, что:
\begin{itemize}
    \item Концепции \textbf{Pointer-Based Security} могут обеспечивать системы с верифицируемыми характеристиками честности
    \item Принципы \textbf{Local Data Regeneration} позволяют поддерживать согласованное состояние без непрерывной передачи данных
    \item Архитектурные комбинации показывают потенциал для возникновения свойств безопасности и снижения зависимостей от инфраструктуры
\end{itemize}

\section{Связь с теоретической основой}

Архитектура исследовательского прототипа соответствует трансформациям, описанным в теоретических работах:

\subsection{Трансформация 1: От передачи данных к синхронному обнаружению}
Состояния игры и решения управляются через скоординированную регенерацию, а не через сетевую передачу. Клиенты используют общие ссылки для поддержания согласованности состояния.

\subsection{Трансформация 2: От хранения к паттернам регенерации}
Архитектура демонстрирует паттерны, в которых необходимые элементы симуляции могут быть регенерированы из минимальных начальных состояний, а не поддерживаться в постоянном хранилище.

\subsection{Трансформация 3: Снижение поверхности атак}
Архитектурные паттерны показывают потенциал для сокращения уязвимых интерфейсов путем минимизации воздействия данных и требований к передаче.

\section{Валидация теоретической парадигмы}

Исследовательский прототип предоставляет экспериментальную поддержку теоретических основ:

\subsection{Поддержка принципов Pointer-Based Security}

Экспериментальные результаты соответствуют теоретическим прогнозам:
\begin{itemize}
    \item \textbf{Сокращение передачи данных} — Архитектурные паттерны минимизируют обмен конфиденциальными данными
    \item \textbf{Координация на основе ссылок} — Открытые координаты обеспечивают синхронизацию состояния
    \item \textbf{Возникающие свойства безопасности} — Архитектурные паттерны обеспечивают неотъемлемые характеристики защиты
    \item \textbf{Минимизация метаданных} — Паттерны обмена показывают сокращение анализируемых метаданных
\end{itemize}

\subsection{Поддержка постулатов Local Data Regeneration}

Экспериментальные доказательства подтверждают теоретические постулаты:

\subsubsection{Постулат 1: Данные как вычисляемое состояние}
Исследование показывает, что состояния игры могут рассматриваться как вычисляемые, а не передаваемые сущности. Состояние $D$ может быть получено через вычисление $F(S, P)$ из ссылок $S$ и $P$.

\subsubsection{Постулат 2: Осуществимость синхронной регенерации}
Экспериментальное подтверждение того, что скоординированные системы могут достигать идентичных состояний $D$ через синхронное применение общих алгоритмов $F$ и ссылок.

\subsubsection{Постулат 3: Коммуникация как координация}
Архитектура демонстрирует коммуникацию в первую очередь как синхронизацию ссылок, а не передачу состояний.

\textbf{Результаты исследований:}
\begin{itemize}
    \item Согласованность состояния из общих ссылок подтверждает Постулат 2
    \item Паттерны доступа к вычислительному состоянию соответствуют Постулату 1
    \item Координация на основе ссылок реализует концепции Постулата 3
    \item Массовая симуляция с минимальным обменом данными подтверждает осуществимость парадигмы
\end{itemize}

\section{Архитектурные паттерны и результаты исследований}

\subsection{Принципы исследования}
\begin{itemize}
    \item \textbf{Детерминированные паттерны:} Согласованные результаты из идентичных входных данных в различных системах
    \item \textbf{Вычисления по требованию:} Состояния вычисляются при необходимости, а не сохраняются постоянно
    \item \textbf{Вселенные на основе ссылок:} Минимальные ссылки определяют сложные состояния систем
\end{itemize}

\subsection{Продемонстрированные исследованием возможности}
\begin{itemize}
    \item \textbf{Процедурная генерация:} Создание контента без постоянного хранения
    \item \textbf{Миры на основе ссылок:} Сложные состояния из минимальных ссылок
    \item \textbf{Динамическое изменение среды:} Множественные конфигурации из модификаций ссылок
    \item \textbf{Симуляция массовых сущностей:} Множество автономных сущностей с минимальной координацией
    \item \textbf{Верифицируемая согласованность:} Математическая проверка целостности состояния
    \item \textbf{Эффективный доступ к состоянию:} Быстрое извлечение и верификация состояния
    \item \textbf{Архитектурная безопасность:} Неотъемлемая защита через паттерны проектирования
    \item \textbf{Сокращенная инфраструктура:} Минимальная зависимость от серверов для основной логики
    \item \textbf{Характеристики масштабируемости:} Линейная производительность с ростом сложности
    \item \textbf{Кроссплатформенная согласованность:} Идентичное поведение в различных системах
    \item \textbf{Детерминированная случайность:} Воспроизводимые случайные последовательности
    \item \textbf{Верифицируемая честность:} Математически доказуемая целостность системы
    \item \textbf{Динамическое изменение системы:} Изменения во время выполнения через обновления ссылок
    \item \textbf{Паттерны без сохранения состояния:} Генерация состояния вместо сохранения
\end{itemize}

\section{Экспериментальные результаты и характеристики производительности}

\subsection{Производительность генерации мира}
\begin{itemize}
    \item \textbf{Тестовая среда:} Конфигурация элементов 1000x1000
    \item \textbf{Производительность генерации:} $\approx$ 0.35 секунды время обработки
    \item \textbf{Характеристики пропускной способности:} $\approx$ 2.8 миллионов элементов/сек
    \item \textbf{Детерминированная согласованность:} Идентичные выходные данные из идентичных ссылок при различных запусках
\end{itemize}

\subsection{Характеристики доступа к состоянию}
Исследовательский прототип демонстрирует ключевое свойство Local Data Regeneration Paradigm: время доступа к состоянию не зависит от его позиционного индекса в огромном пространстве состояний. Измерения показывают согласованное время доступа для диапазона позиций от 1 до $10^{100}$:

\begin{center}
\begin{tabular}{|l|r|r|}
\hline
\textbf{Позиция} & \textbf{Время (сек)} & \textbf{Верифицировано} \\
\hline
1 & 0.00010562 & Да \\
1K & 0.00002575 & Да \\
1M & 0.00001550 & Да \\
1B & 0.00001359 & Да \\
1T & 0.00001502 & Да \\
1Q & 0.00001478 & Да \\
$10^{20}$ & 0.00001502 & Да \\
$10^{100}$ & 0.00002027 & Да \\
\hline
\end{tabular}
\end{center}

Измеренное время доступа ($\approx$ 0.000015-0.000020 секунд) остается согласованным для всех протестированных позиций, демонстрируя, что производительность извлечения состояния фактически не зависит от индекса позиции состояния.

\subsection{Производительность принятия решений сущностями}

Симуляция сущностей показывает линейные характеристики масштабирования. В отличие от систем, зависящих от сети, где дополнительные сущности влияют на производительность, данная архитектура поддерживает согласованную производительность на сущность.

\begin{center}
\begin{tabular}{|l|r|r|r|r|}
\hline
\textbf{Режим} & \textbf{Сущности} & \textbf{Операции} & \textbf{Длительность (сек)} & \textbf{Скорость (оп/сек)} \\
\hline
Производ. & 100 & 100,000 & 0.037308 & 2,680,375 \\
\hline
Производ. & 1,000 & 1,000,000 & 0.337940 & 2,959,105 \\
\hline
Вериф. & 100 & 100,000 & 0.205825 & 485,849 \\
\hline
Вериф. & 1,000 & 1,000,000 & 2.019662 & 495,132 \\
\hline
\end{tabular}
\end{center}

\begin{itemize}
    \item \textbf{Режим производительности:} До 2.9 миллионов операций/сек (совокупно)
    \item \textbf{Верифицированный режим:} До 495,000 верифицированных операций/сек (совокупно)
    \item \textbf{Характеристики масштабирования:} Линейное увеличение времени с количеством сущностей, указывающее на сложность O(n)
    \item \textbf{Влияние верификации:} Криптографическая верификация показывает примерно 6-кратное влияние на производительность
\end{itemize}

\subsection{Возможность верификации решений}
Исследование продемонстрировало возможность верификации операций за постоянное время, включая верификацию 1,000,000,000-й операции без последовательной обработки.

\section{Последствия исследований и потенциальное влияние}

Экспериментальные результаты предполагают значительные потенциальные последствия:
\begin{itemize}
    \item \textbf{Эффективность инфраструктуры:} Архитектурные паттерны предполагают потенциальное снижение зависимости от серверов
    \item \textbf{Характеристики безопасности:} Паттерны проектирования показывают неотъемлемые свойства защиты
    \item \textbf{Эффективность разработки:} Упрощенные требования к синхронизации
    \item \textbf{Генерация контента:} Возможности динамического создания контента
    \item \textbf{Согласованность платформ:} Единообразие поведения на различных платформах
\end{itemize}

\section{Заключение}

Данное исследование предоставляет экспериментальные доказательства, подтверждающие осуществимость теоретических парадигм, представленных в \cite{suvorov2025pointer} и \cite{suvorov2025local}. Исследовательский прототип, реализованный как Deterministic Game Engine, служащий модельной средой, демонстрирует, что:

\textbf{Архитектурные сдвиги от парадигм передачи данных к парадигмам регенерации показывают потенциал для создания систем с улучшенными характеристиками производительности, безопасности и масштабируемости.}

Экспериментальные результаты поддерживают обе теоретические основы: Pointer-Based Security через паттерны сокращения передачи данных и Local Data Regeneration через реализацию основных постулатов. Демонстрация времени доступа к состоянию, не зависящего от индекса позиции, и линейного масштабирования во время симуляции массовых сущностей предоставляет конкретную валидацию преимуществ парадигм над традиционными сетезависимыми архитектурами.

Будущие исследования будут исследовать применение этих архитектурных паттернов в более широких областях, включая распределенные симуляции, системы IoT и верифицируемые вычислительные среды.

\section*{Благодарности}
Автор благодарит исследовательское сообщество за ценные обсуждения теоретических основ.

\begin{thebibliography}{9}
\bibitem{suvorov2025pointer} Suvorov, A. (2025). The Pointer-Based Security Paradigm: Architectural Shift from Data Protection to Data Non-Existence. Zenodo. \url{https://doi.org/10.5281/zenodo.17204738}
\bibitem{suvorov2025local} Suvorov, A. (2025). The Local Data Regeneration Paradigm: Ontological Shift from Data Transmission to Synchronous State Discovery. Zenodo. \url{https://doi.org/10.5281/zenodo.17264327}
\end{thebibliography}

\end{document}